\section{从电流积分到 $a_E/a_M$}
%=============================================================================

本节是推导的核心：将标量波动方程的解代入 Eqs.(3)--(4)，得到 $a_E/a_M$ 以 $\mathbf{J}_S$ 表示的积分式。

\subsection{标量波动方程的 Green 函数解}

对于形如 $(\nabla^2 + k^2)\psi(\mathbf{r}) = -f(\mathbf{r})$ 的标量 Helmholtz 方程，
出射波边条件的解为：
$$
\psi(\mathbf{r}) = \int G(\mathbf{r},\mathbf{r}')\,f(\mathbf{r}')\,d^3r'
$$
其中出射 Green 函数为：
$$
G(\mathbf{r},\mathbf{r}') = \frac{e^{ik|\mathbf{r}-\mathbf{r}'|}}{4\pi|\mathbf{r}-\mathbf{r}'|}
$$

利用球谐函数的加法公式展开：
$$
G(\mathbf{r},\mathbf{r}') = ik\sum_{l=1}^\infty\sum_{m=-l}^l j_l(kr_<)\,h_l^{(1)}(kr_>)\,
  Y_{lm}(\theta,\phi)\,Y_{lm}^*(\theta',\phi')
$$
其中 $r_< = \min(r,r')$，$r_> = \max(r,r')$。

\subsection{$a_M$ 的电流积分形式}

我们需要将 $\mathbf{r}\cdot\mathbf{H}_s$ 代入 $a_M$ 的表达式。从 \S4 我们知道 $a_M$ 可以通过
$\mathbf{H}_s$ 表达。利用对偶性，从 Eq.(2) 出发，用 $\hat{\mathbf{r}}\cdot$ 提取 $a_M$：

由 §3 的 Grahn Eq.(2)：$\mathbf{H}_s$ 的展开式中 $a_M$ 出现在 $\nabla\times$ 项（含径向分量），
$a_E$ 出现在 $\mathbf{X}_{lm}$ 项（纯切向）。因此提取 $a_M$ 可以用 $\mathbf{H}_s$ 的径向分量做投影
（Grahn 的 Eq.(5) 对偶路径）。

从 Eq.(2) 的径向分量：
$$
\hat{\mathbf{r}}\cdot\mathbf{H}_s = \frac{E_0}{\eta}\sum_{lm} i^{l-1}[\pi(2l+1)]^{1/2}
  a_M(l,m)\,\frac{\sqrt{l(l+1)}}{kr}\,h_l^{(1)}(kr)\,Y_{lm}
$$

因此：
$$
a_M(l,m) = \frac{(-i)^{l}kr\,\eta}{h_l^{(1)}(kr)E_0[\pi(2l+1)l(l+1)]^{1/2}}
  \int Y_{lm}^*\,\hat{\mathbf{r}}\cdot\mathbf{H}_s\,d\Omega
$$
（这是从 $\mathbf{H}_s$ 径向投影的对偶路径，与 Eq.(3) 的 $\mathbf{E}_s$ 投影不同；
中间态相位 $(-i)^l$ 与随后 Green 展开的 $ik$ 因子中的 $i$ 合并为 $(-i)^{l-1}$，
最终与 Grahn Eq.(14) 一致。原文 Eq.(5) 的中间相位需对照 Grahn 原文确认。）

$\blacktriangleright$ 现在代入 $\mathbf{r}\cdot\mathbf{H}_s$ 的 Green 函数解 (Eq.(12) 的解)。

Eq.(12) 写为标准 Helmholtz 形式：$(\nabla^2+k^2)(\mathbf{r}\cdot\mathbf{H}_s) = -\mathbf{r}\cdot(\nabla\times\mathbf{J}_S)$。

在散射体外部（取球面半径 $r$ 大于所有源点 $r'$），$r_> = r$，$r_< = r'$：
$$
\mathbf{r}\cdot\mathbf{H}_s(\mathbf{r}) = ik\int \sum_{lm} j_l(kr')\,h_l^{(1)}(kr)\,
  Y_{lm}(\theta,\phi)\,Y_{lm}^*(\theta',\phi')\,
  \bigl[\mathbf{r}'\cdot(\nabla'\times\mathbf{J}_S(\mathbf{r}'))\bigr]\,d^3r'
$$

因此：
$$
\hat{\mathbf{r}}\cdot\mathbf{H}_s = \frac{ik}{r}\int \sum_{lm} j_l(kr')\,h_l^{(1)}(kr)\,
  Y_{lm}(\theta,\phi)\,Y_{lm}^*(\theta',\phi')\,
  \bigl[\mathbf{r}'\cdot(\nabla'\times\mathbf{J}_S)\bigr]\,d^3r'
$$

乘以 $Y_{l'm'}^*$ 并在球面积分：
$$
\int Y_{l'm'}^*\,\hat{\mathbf{r}}\cdot\mathbf{H}_s\,d\Omega
= \frac{ik}{r} h_l^{(1)}(kr)\,
  \sum_{lm} \delta_{ll'}\delta_{mm'}\,
  \int j_l(kr')\,Y_{lm}^*(\theta',\phi')\,
  \bigl[\mathbf{r}'\cdot(\nabla'\times\mathbf{J}_S)\bigr]\,d^3r'
$$
$$
= \frac{ik}{r} j_{l'}(kr')\,h_{l'}^{(1)}(kr)\,
  \int Y_{l'm'}^*(\theta',\phi')\,
  \bigl[\mathbf{r}'\cdot(\nabla'\times\mathbf{J}_S)\bigr]\,d^3r'
$$
（注意 $j_{l'}(kr')$ 依赖积分变量 $r'$，故移入体积分内。）

代回 $a_M$ 表达式（注意 Green 展开的 $ik$ 因子中的 $i$ 与 $(-i)^l$ 合并：$(-i)^l\cdot i=(-i)^{l-1}$）：
$$
a_M(l,m) = \frac{(-i)^{l}kr\,\eta}{h_l^{(1)}(kr)E_0[\pi(2l+1)l(l+1)]^{1/2}}
  \cdot\frac{ik}{r} h_l^{(1)}(kr)
  \int j_l(kr')\,Y_{lm}^*(\theta',\phi')\,
  \mathbf{r}'\cdot(\nabla'\times\mathbf{J}_S)\,d^3r'
$$

$h_l^{(1)}(kr)$ 和 $r$ 消去，$(-i)^l\cdot i=(-i)^{l-1}$ 后，整理得到：
\begin{equation}
  \boxed{\;
    a_M(l,m) = \frac{(-i)^{l-1}k^2\eta}{E_0[\pi(2l+1)l(l+1)]^{1/2}}
    \int Y_{lm}^*(\theta,\phi)\,j_l(kr)\,
    \mathbf{r}\cdot\bigl[\nabla\times\mathbf{J}_S(\mathbf{r})\bigr]\,d^3r
    \;}
  \label{eq:aM_J}
\end{equation}

这就是 Grahn 的 Eq.(14)。\textbf{磁多极系数来自 $\mathbf{J}_S$ 的旋度}。

\subsection{$a_E$ 的电流积分形式}

类似地，利用 Eq.(11) 的解代入 Eq.(3)。

Eq.(11) 写为：$(\nabla^2+k^2)(\mathbf{r}\cdot\mathbf{E}_s) = S_E(\mathbf{r})$，其中
$$
S_E(\mathbf{r}) = \frac{2\rho_S + r\partial_r\rho_S}{\epsilon_0\epsilon_{r,d}} - i\omega\mu_0\,\mathbf{r}\cdot\mathbf{J}_S
$$

利用连续性方程 $\nabla\cdot\mathbf{J}_S = i\omega\rho_S$，可以将 $\rho_S$ 项用 $\nabla\cdot\mathbf{J}_S$ 表示：
$$
\frac{2\rho_S + r\partial_r\rho_S}{\epsilon_0\epsilon_{r,d}} = \frac{1}{i\omega\epsilon_0\epsilon_{r,d}}\bigl(2\nabla\cdot\mathbf{J}_S + r\partial_r(\nabla\cdot\mathbf{J}_S)\bigr)
$$

其中 $(2+r\partial_r)(\nabla\cdot\mathbf{J}_S)$ 正是 Eq.(13) 的第二项。经过与 $a_M$ 类似的 Green 函数代入 + 正交投影步骤：

\begin{equation}
  \boxed{\;
    a_E(l,m) = \frac{(-i)^{l-1}k\eta}{E_0[\pi(2l+1)l(l+1)]^{1/2}}
    \int Y_{lm}^*(\theta,\phi)\,j_l(kr)\Bigl[
      k^2\mathbf{r}\cdot\mathbf{J}_S + \left(2 + r\frac{d}{dr}\right)(\nabla\cdot\mathbf{J}_S)
    \Bigr] d^3r
    \;}
  \label{eq:aE_J}
\end{equation}

这就是 Grahn 的 Eq.(13)。\textbf{电多极系数来自 $\mathbf{J}_S$ 的散度}。

\begin{stepbox}{为什么 $a_E$ 含 $\nabla\cdot\mathbf{J}_S$ 而 $a_M$ 含 $\nabla\times\mathbf{J}_S$？}
这是由源在 Maxwell 方程组中的角色决定的：
\begin{itemize}
  \item 电多极与电荷分布（$\nabla\cdot\mathbf{J}_S\sim\rho_S$）关联——电场有散度源
  \item 磁多极与环流分布（$\nabla\times\mathbf{J}_S$）关联——磁场有旋度源
\end{itemize}
这一区别在下一节分部积分后仍会保留。
\end{stepbox}

\subsection{为什么 Eqs.(13)--(14) 数值计算不便？}

\begin{noteworthy}
虽然 Eqs.(13)--(14) 在理论上是干净的，但\textbf{数值实现存在严重问题}：
\begin{enumerate}
  \item \textbf{需要计算 $\nabla\cdot\mathbf{J}_S$ 和 $\nabla\times\mathbf{J}_S$}：
    在 FDTD 或 FEM 网格上，$\mathbf{J}_S$ 是离散的，数值求导会放大离散误差，
    尤其是对高阶多极（$l$ 大时角向变化快）。
   \item \textbf{$\left(2+r\frac{d}{dr}\right)$ 算子}作用在 $\nabla\cdot\mathbf{J}_S$ 上：
     这要求对散度场再做径向微分——对网格数据而言精度损失严重。
  \item \textbf{不可分离性}：$\nabla\cdot\mathbf{J}_S$ 和 $\nabla\times\mathbf{J}_S$ 混叠了不同粒子的贡献，
    不利于逐个粒子分析。
\end{enumerate}
\textbf{解决方案}：通过分部积分将导数从 $\mathbf{J}_S$ 转移到已知解析函数（球谐、Bessel 函数）上。
\end{noteworthy}

\begin{chaptersummary}

\textbf{§6 章节总结}

本章将标量波动方程的解代入 §4 的投影公式，得到 $a_E/a_M$ 的电流积分形式。

\textbf{Green 函数法：} Helmholtz 方程 $(\nabla^2+k^2)\psi = -f$ 的出射波解为
$$
\psi(\mathbf{r}) = \int G(\mathbf{r},\mathbf{r}')\,f(\mathbf{r}')\,d^3r',\quad
G = \frac{e^{ik|\mathbf{r}-\mathbf{r}'|}}{4\pi|\mathbf{r}-\mathbf{r}'|}
$$
在球坐标中展开为球 Bessel 函数和球谐的谱形式。

\textbf{a$_M$ 的电流形式（Eq.(14)）——来自 $\mathbf{r}\cdot\mathbf{H}_s$ 的波动方程：}
$$
\boxed{\;
a_M(l,m) = \frac{(-i)^{l-1}k^2\eta}{E_0[\pi(2l+1)l(l+1)]^{1/2}}
\int Y_{lm}^*\,j_l(kr)\;\mathbf{r}\cdot(\nabla\times\mathbf{J}_S)\;d^3r
\;}
$$

\textbf{a$_E$ 的电流形式（Eq.(13)）——来自 $\mathbf{r}\cdot\mathbf{E}_s$ 的波动方程：}
$$
\boxed{\;
a_E(l,m) = \frac{(-i)^{l-1}k\eta}{E_0[\pi(2l+1)l(l+1)]^{1/2}}
\int Y_{lm}^*\,j_l(kr)\Bigl[
      k^2\mathbf{r}\cdot\mathbf{J}_S + \left(2 + r\frac{d}{dr}\right)(\nabla\cdot\mathbf{J}_S)
\Bigr] d^3r
\;}
$$

\textbf{电/磁多极对 $\mathbf{J}_S$ 运算的差异：}
\begin{itemize}
  \item $a_M$ 需要 $\nabla\times\mathbf{J}_S$（旋度，来自磁场的波动方程 Eq.(12)）。
  \item $a_E$ 需要 $\nabla\cdot\mathbf{J}_S$ 及其径向导数（散度，来自电场的波动方程 Eq.(11) 和连续性方程）。
\end{itemize}
这反映了电磁对偶：旋度→磁多极、散度→电多极。

\textbf{Eqs.(13)--(14) 的理论意义：} 它们将 $a_E/a_M$ 与 $\mathbf{J}_S$ 直接关联，
不再需要包围球面上的 $\mathbf{E}_s$ 值，而是只需在散射体体积内计算 $\mathbf{J}_S$ 的积分。
但数值实现仍有困难——下一节解决。

\end{chaptersummary}

\newpage

%=============================================================================
